% ============================================================
% SLIDE 1: Reward-Berechnung
% ============================================================
\begin{frame}{Methoden -- Reward-Berechnung}
\begin{columns}[T]
\begin{column}{0.48\textwidth}
  \begin{block}{Dichte Belohnungen (pro Schritt)}
    \begin{itemize}
      \item \textbf{Distanz-Penalty}\\
            $r_\text{dist} = -5{,}0 \cdot \Delta t \cdot \|\mathbf{p}_\text{rel}\|$
      \item \textbf{Zeit-Penalty}\\
            $r_\text{time} = -0{,}5 \cdot \Delta t$
    \end{itemize}
    \vspace{0.2em}
    {\small Skalierung mit $\Delta t = 0{,}01\,$s macht Rewards zeitschrittunabhängig.}
  \end{block}
\end{column}
\begin{column}{0.48\textwidth}
  \begin{exampleblock}{Terminale Belohnungen (einmalig)}
    \begin{itemize}
      \item \textbf{Crash}: $r_\text{crash} = -100$
      \item \textbf{Erfolg}: $r_\text{success} = +200$
    \end{itemize}
    \vspace{0.2em}
    {\small Nicht mit $\Delta t$ skaliert -- einmalige Ereignisse.}
  \end{exampleblock}
  \begin{alertblock}{Gesamtreward}
    $R = \sum_t (r_\text{dist} + r_\text{time}) + r_\text{terminal}$
  \end{alertblock}
\end{column}
\end{columns}
\end{frame}


% ============================================================
% SLIDE 2: Netzwerk-Architektur
% ============================================================
\begin{frame}{Methoden -- Netzwerk-Architektur}
\begin{columns}[T]
\begin{column}{0.48\textwidth}
  \begin{block}{Actor-Critic (PPO)}
    \begin{itemize}
      \item \textbf{Actor}: MLP $[128, 128]$
      \item \textbf{Critic}: MLP $[128, 128]$
      \item Aktivierung: \texttt{tanh}
      \item Initiales Rauschen: $\sigma_0 = 1{,}0$
    \end{itemize}
  \end{block}
  \vspace{0.2cm}
  \begin{block}{Eingabe / Ausgabe}
    \begin{itemize}
      \item Eingabe: 17 Observationen
      \item Ausgabe: 4 Aktionen $\in [-1, 1]$
    \end{itemize}
  \end{block}
\end{column}
\begin{column}{0.48\textwidth}
  \begin{exampleblock}{Datenfluss}
    \small
    \textbf{Obs} (17) $\xrightarrow{\text{Actor}}$ \textbf{Aktion} (4)\\
    \vspace{0.3em}
    $\downarrow$ \textit{Action Scaling}\\
    \vspace{0.3em}
    \textbf{PID-Zielposition} $(x, y, z, \psi)$\\
    \vspace{0.3em}
    $\downarrow$ \textit{Kaskadierender PID}\\
    \vspace{0.3em}
    \textbf{Motor-RPMs} (4 Propeller)\\
    \vspace{0.3em}
    $\downarrow$ \textit{Genesis Physik}\\
    \vspace{0.3em}
    \textbf{Neuer Zustand}
  \end{exampleblock}
\end{column}
\end{columns}
\end{frame}


% ============================================================
% SLIDE 3: Observationen und Normalisierung
% ============================================================
\begin{frame}{Methoden -- Observationsraum}
\begin{columns}[T]
\begin{column}{0.55\textwidth}
  \begin{block}{Observationsvektor ($\dim = 17$)}
    \renewcommand{\arraystretch}{1.2}
    \begin{tabular}{lll}
      \hline
      \textbf{Größe} & \textbf{Dim} & \textbf{Skalierung} \\
      \hline
      Relative Position & 3 & $\times \tfrac{1}{15}$, clip $[-1,1]$ \\
      Quaternion        & 4 & -- \\
      Lineare Geschw.   & 3 & $\times \tfrac{1}{5}$, clip $[-1,1]$ \\
      Winkelgeschw.     & 3 & $\times \tfrac{1}{\pi}$, clip $[-1,1]$ \\
      Letzte Aktionen   & 4 & -- \\
      \hline
    \end{tabular}
  \end{block}
\end{column}
\begin{column}{0.42\textwidth}
  \begin{exampleblock}{Warum Normalisierung?}
    \begin{itemize}
      \item Werte $\in [-1, 1]$ stabilisieren das Training
      \item Skalierung richtet sich nach physikalisch sinnvollen Maxima
      \item Relative Position statt absoluter: generalisiert besser
    \end{itemize}
  \end{exampleblock}
\end{column}
\end{columns}
\end{frame}


% ============================================================
% SLIDE 4: Hyperparameter
% ============================================================
\begin{frame}{Methoden -- Hyperparameter}
\begin{columns}[T]
\begin{column}{0.48\textwidth}
  \begin{block}{PPO / Lernparameter}
    \renewcommand{\arraystretch}{1.2}
    \begin{tabular}{ll}
      \hline
      \textbf{Parameter}       & \textbf{Wert} \\
      \hline
      Discount $\gamma$        & $0{,}99$ \\
      GAE $\lambda$            & $0{,}95$ \\
      Lernrate                 & $3 \times 10^{-4}$ \\
      Clipping $\epsilon$      & $0{,}2$ \\
      Epochen / Update         & $5$ \\
      Mini-Batches             & $4$ \\
      Schritte pro Env         & $100$ \\
      \hline
    \end{tabular}
  \end{block}
\end{column}
\begin{column}{0.48\textwidth}
  \begin{exampleblock}{Reward-Skalierung}
    \begin{tabular}{ll}
      \hline
      \textbf{Reward} & \textbf{Skala} \\
      \hline
      Distanz  & $-5{,}0$ \\
      Zeit     & $-0{,}5$ \\
      Crash    & $-100{,}0$ \\
      Erfolg   & $+200{,}0$ \\
      \hline
    \end{tabular}
  \end{exampleblock}
  \begin{exampleblock}{Aktionsskalierung}
    $\Delta x, \Delta y, \Delta z \in [-1, +1]\,$m\\
    $\psi_\text{target} = a_\psi \cdot 180°$
  \end{exampleblock}
\end{column}
\end{columns}
\end{frame}


% ============================================================
% SLIDE 5: Environment-Architektur und Parameter
% ============================================================
\begin{frame}{Methoden -- Environment-Parameter}
\begin{columns}[T]
\begin{column}{0.48\textwidth}
  \begin{block}{Spawn \& Zielbereich}
    \begin{itemize}
      \item \textbf{Drohne}: $x, y \in [-5, +5]\,$m\\
            Höhe: $z = 10\,$m (fest)
      \item \textbf{Ziel (PPO)}: fest bei $(3, 3, 1)\,$m
      \item \textbf{Episodendauer}: $30\,$s\\
            ($= 3000$ Schritte bei $100\,$Hz)
      \item \textbf{Parallelität}: 4096 Umgebungen
    \end{itemize}
  \end{block}
  \begin{block}{Curriculum}
    \begin{itemize}
      \item Erste $\mathbf{20\,000}$ Schritte:\\
            Ziel $\leq 1\,$m von Drohne entfernt
      \item Danach: volles Zielraster
    \end{itemize}
  \end{block}
\end{column}
\begin{column}{0.48\textwidth}
  \begin{alertblock}{Abbruchbedingungen}
    \begin{itemize}
      \item Höhe $< 0{,}2\,$m (Boden)
      \item Neigung $> 60°$ (Roll oder Pitch)
      \item Abstand zum Ziel $> 50\,$m
      \item Zeitlimit überschritten ($30\,$s)
    \end{itemize}
  \end{alertblock}
  \begin{exampleblock}{Erfolgsbedingung}
    $30$ konsekutive Schritte mit:\\
    $\|\mathbf{p}_\text{rel}\| < 0{,}3\,$m \textbf{und}\\
    $\|\mathbf{v}\| < 0{,}3\,$m/s
    \vspace{0.2em}\\
    {\small ($\hat{=}\ 0{,}3\,$s Schwebezeit am Ziel)}
  \end{exampleblock}
\end{column}
\end{columns}
\end{frame}


% ============================================================
% SLIDE 6: Abschlussfolie
% ============================================================
{
\setbeamertemplate{footline}{}
\begin{frame}[plain]
\begin{tikzpicture}[remember picture, overlay]
  \fill[svgPurple] (current page.north west) rectangle (current page.south east);
  \fill[svgPurpleDark, opacity=0.4] (current page.south west) rectangle
    ([yshift=3cm]current page.south east);
  \fill[white, opacity=0.04] ([xshift=-2cm, yshift=-1cm]current page.north east) circle (4cm);
  \fill[white, opacity=0.03] ([xshift=1cm, yshift=1cm]current page.south west) circle (5cm);
  \draw[svgOrange, line width=2pt]
    ([xshift=2cm, yshift=-4.8cm]current page.north west) --
    ([xshift=14cm, yshift=-4.8cm]current page.north west);
\end{tikzpicture}
\vspace{1.5cm}
\begin{center}
  {\Large\bfseries\color{white}Vielen Dank}\\[0.3em]
  {\normalsize\color{svgPurpleLight}für Ihre Aufmerksamkeit}\\[1em]
  {\color{svgOrangeLight}\rule{4cm}{0.5pt}}\\[0.8em]
  {\normalsize\color{svgPurpleLight}Fragen?}\\[1.5em]
  {\small\color{white}Fritz Körner}
\end{center}
\end{frame}
}
